%!TEX TS-program = xelatex
%!TEX encoding = UTF-8
% 经济研究（Economic Research Journal）格式模板
% 参考来源：http://www.erjw.cn/
% 适用于经济金融领域实证论文
% CTeX 中文格式，使用XeLaTeX编译

\PassOptionsToPackage{quiet}{xeCJK}
\documentclass[lang=cn, alanguage={chinese}, aaffiliation={中国}, aoptions={authoryear}]{ctexart}

% ── 文档类型选项 ────────────────────────────────────────────────────────────
% 本模板采用单栏排版，如需双栏可取消下面注释：
% \documentclass[lang=cn, twocolumn]{ctexart}

% ── 必要宏包 ────────────────────────────────────────────────────────────────
\usepackage[
    a4paper,
    top=2.5cm, bottom=2.5cm,
    left=3.0cm, right=2.5cm,
    headheight=14.5pt,
    heightrounded,
]{geometry}

\usepackage{fontspec}           % XeLaTeX 字体配置
\usepackage{xeCJK}              % 中日韩文字处理
\setmainfont{Times New Roman}   % 英文正文
\setsansfont{Arial}
\setmonofont{Courier New}
\xeCJKsetcharclass("0, "7F) { } % 忽略 ASCII 控制字符
\xeCJKsetup{
    CJKspace=true,
    xCJKecglue={},
    CheckFullChar=true,
}

\usepackage{mathtools}          % 数学公式（需 amsmath）
\usepackage{booktabs}           % 表格线（\toprule/midrule/bottomrule）
\usepackage{threeparttable}     % 表格注释
\usepackage{longtable}          % 跨页长表格
\usepackage{graphicx}           % 插图
\graphicspath{{figures/}}       % 图表目录
\usepackage{subcaption}         % 子图 subfigure 替代
\usepackage{hyperref}           % 超链接与 PDF 元数据
\hypersetup{
    colorlinks=true,
    linkcolor=blue,
    citecolor=blue,
    urlcolor=blue,
    pdftitle={经济研究格式模板},
    pdfauthor={作者姓名},
}

\usepackage{natbib}             % 参考文献（GB/T 7714-2015 顺序编码制）
\bibliographystyle{gbt7714-numerical}
% 可选风格：gbt7714-author-year（著者-出版年制）
% 需配合 biber 使用：xelatex → biber → xelatex × 2

\usepackage{setspace}           % 行距设置
\usepackage{enumitem}           % 列表设置
\usepackage{indentfirst}        % 首段缩进
\usepackage{titlesec}           % 章节标题格式

% ── 页面设置 ────────────────────────────────────────────────────────────────
% 经济研究要求：小四号（12pt）字体，1.5 倍行距
\setstretch{1.5}
\setlength{\parindent}{2em}     % 首行缩进 2 个中文字符
\setlength{\abovecaptionskip}{6pt}
\setlength{\belowcaptionskip}{6pt}
\linespread{1.5}

% ── 章节标题格式 ────────────────────────────────────────────────────────────
\titleformat*{\section}{\fontsize{13pt}{1.5em}\bfseries\center}
\titleformat*{\subsection}{\fontsize{12pt}{1.5em}\bfseries}
\titleformat*{\subsubsection}{\fontsize{12pt}{1.5em}\rmfamily}

% ── 页眉页脚 ────────────────────────────────────────────────────────────────
\usepackage{fancyhdr}
\pagestyle{fancy}
\fancyhf{}
\fancyhead[C]{\footnotesize 中国工业经济}
\fancyfoot[C]{\thepage}
% 双栏模式下使用：
% \fancyhead[LE,RO]{\thepage}
% \fancyhead[LO]{\leftmark}
% \fancyhead[RE]{\rightmark}

% ── 中文破折号与省略号 ──────────────────────────────────────────────────────
\ExplSyntaxOn
\DeclareDocumentCommand \zhdash}{ }{
    \unskip\ignorespaces\textemdash\ignorespaces}
\ExplSyntaxOff

% ── 标题与作者信息 ──────────────────────────────────────────────────────────
\title{论文标题（中文）\thanks{基金项目（编号）。}}
\author{作者甲\thanks{作者简介：...，E-mail: ...} \quad 作者乙\thanks{...}}
\date{\today}

% ── 摘要环境 ────────────────────────────────────────────────────────────────
\newenvironment{cnabstract}{%
    \par\textbf{摘要：}
    \@tempswafalse
}{\par\vspace{\baselineskip}\par\textbf{关键词：}...}

\newenvironment{enabstract}{%
    \par\begin{minipage}{\textwidth}\textbf{Abstract: }
    }{%
    \end{minipage}\par\vspace{\baselineskip}\par\textbf{Keywords: }...}

% ═══════════════════════════════════════════════════════════════════════════
%                            正文内容
% ═══════════════════════════════════════════════════════════════════════════

\begin{document}

% 消除首页页眉
\thispagestyle{empty}

% 标题（中英文）
\maketitle
\thispagestyle{fancy}

% ── 中文摘要 ────────────────────────────────────────────────────────────────
\begin{cnabstract}
    本文基于\dots 研究了\dots。实证结果表明\dots。
    本文的研究贡献在于：第一，\dots；第二，\dots。
\end{cnabstract}

\textbf{关键词：}关键词1；关键词2；关键词3；关键词4

% ── 英文摘要（可选） ────────────────────────────────────────────────────────
\begin{enabstract}
    This paper investigates \dots using \dots data.
    The results show that \dots. Our contributions are twofold: first, \dots;
    second, \dots.
\end{enabstract}

\textbf{Keywords:} keyword1; keyword2; keyword3; keyword4

% ── 一、引言 ───────────────────────────────────────────────────────────────
\section{引言}
\label{sec:intro}

本文研究的问题是\dots。

首先，\dots。其次，\dots。

本文的结构安排如下：第二部分为文献综述；第三部分为研究设计；第四部分为实证结果与分析；第五部分为进一步分析；第六部分为结论与启示。

% ── 二、文献综述 ─────────────────────────────────────────────────────────
\section{文献综述与研究假说}
\label{sec:lit}

\subsection{文献综述}
\label{sec:lit:review}

关于\dots的研究已有丰富的文献，主要分为三类：

第一，\dots（\citealp{xxx}）。\citet{yyy}发现\dots。

第二，\dots（\citealp{zzz}）。\citet{www}利用\dots方法，证实了\dots。

第三，\dots（\citealp{vvv}）。

\subsection{研究假说}
\label{sec:hypothesis}

基于上述分析，本文提出如下假说：

\textbf{假说 H1：}在其他条件不变的情况下，\dots对\dots有显著的正向影响。

\textbf{假说 H2：}\dots对\dots的边际效应存在异质性，表现为\dots。

% ── 三、研究设计 ──────────────────────────────────────────────────────────
\section{研究设计}
\label{sec:design}

\subsection{样本选择与数据来源}
\label{sec:data}

本文以\dots为研究样本，时间跨度为\dots年。数据来源包括：
（1）CSMAR 数据库；（2）Wind 数据库；（3）国家统计局。

本文对原始数据做如下处理：第一，\dots；第二，\dots。

\subsection{变量定义}
\label{sec:var}

\textbf{被解释变量（$Y$）：}本文采用\dots衡量\dots。

\textbf{核心解释变量（$X$）：}本文的核心解释变量为\dots。

\textbf{控制变量（Controls）：}参考\citet{xxx}，本文控制了以下变量：
公司规模（Size）、资产负债率（Lev）、盈利能力（ROA）、企业年龄（Age）等。

表~\ref{tab:var_def}~汇报了主要变量的定义。

\begin{table}[htbp]
    \centering
    \caption{主要变量定义}
    \label{tab:var_def}
    \begin{threeparttable}
        \begin{tabular}{p{2.5cm}p{2.5cm}p{5cm}}
            \toprule
            \textbf{变量名称} & \textbf{变量符号} & \textbf{变量定义} \\
            \midrule
            被解释变量 & $Y$ & \dots \\
            核心解释变量 & $X$ & \dots \\
            控制变量 & Size & \dots \\
            & Lev & \dots \\
            & ROA & \dots \\
            \bottomrule
        \end{tabular}
        \begin{tablenotes}\footnotesize
            \item \textbf{数据来源：}CSMAR 数据库。
            \item 缩尾处理：所有连续变量在 1\% 和 99\% 分位数进行 Winsorize 处理。
        \end{tablenotes}
    \end{threeparttable}
\end{table}

\subsection{模型设定}
\label{sec:model}

为检验假说 H1，本文构建如下双向固定效应模型：

\begin{equation}
    \label{eq:baseline}
    Y_{it} = \alpha + \beta X_{it} + \gamma Controls_{it} + \mu_i + \lambda_t + \varepsilon_{it},
\end{equation}

其中，$i$ 表示企业，$t$ 表示年份。$\mu_i$ 和 $\lambda_t$ 分别为企业和年份固定效应，$\varepsilon_{it}$ 为随机误差项。

% ── 四、实证结果 ──────────────────────────────────────────────────────────
\section{实证结果与分析}
\label{sec:results}

\subsection{描述性统计}
\label{sec:desc}

表~\ref{tab:desc_stat}~报告了主要变量的描述性统计。

\begin{table}[htbp]
    \centering
    \caption{描述性统计}
    \label{tab:desc_stat}
    \begin{threeparttable}
        \begin{tabular}{lcccccc}
            \toprule
            \textbf{变量} & \textbf{N} & \textbf{Mean} & \textbf{SD} & \textbf{Min} & \textbf{Median} & \textbf{Max} \\
            \midrule
            $Y$ & 8,432 & 0.051 & 0.108 & -0.231 & 0.039 & 0.382 \\
            $X$ & 8,432 & 0.312 & 0.463 & 0.000 & 0.000 & 1.000 \\
            Size & 8,432 & 22.14 & 1.326 & 19.76 & 21.95 & 25.89 \\
            Lev & 8,432 & 0.421 & 0.198 & 0.047 & 0.410 & 0.853 \\
            \bottomrule
        \end{tabular}
        \begin{tablenotes}\footnotesize
            \item 注：所有连续变量均经过 1\% 和 99\% 分位数缩尾处理。
        \end{tablenotes}
    \end{threeparttable}
\end{table}

\subsection{基准回归}
\label{sec:baseline}

表~\ref{tab:baseline}~报告了基准回归结果。

\begin{table}[htbp]
    \centering
    \caption{基准回归结果}
    \label{tab:baseline}
    \begin{threeparttable}
        \begin{tabular}{lccc}
            \toprule
            \textbf{} & \textbf{(1)} & \textbf{(2)} & \textbf{(3)} \\
            \midrule
            $X$ & 0.023\*** & 0.018\*** & 0.015\** \\
              & (3.42) & (2.89) & (2.21) \\
            Size & & -0.002 & -0.003 \\
              & & (-0.54) & (-0.71) \\
            \bottomrule
            \textbf{固定效应} & \textbf{否} & \textbf{否} & \textbf{是} \\
            \bottomrule
        \end{tabular}
        \begin{tablenotes}\footnotesize
            \item 注：括号内为 $t$ 值；\* $p<0.1$，\** $p<0.05$，\*** $p<0.01$。
            \item 企业层面聚类标准误，*** $p<0.01$, ** $p<0.05$, * $p<0.10$。
        \end{tablenotes}
    \end{threeparttable}
\end{table}

% ── 五、进一步分析 ────────────────────────────────────────────────────────
\section{进一步分析}
\label{sec:additional}

\subsection{异质性分析}
\label{sec:hetero}

本文从企业类型、地区等维度进行异质性分析，结果报告于表~\ref{tab:hetero}。

\begin{figure}[htbp]
    \centering
    \caption{平行趋势检验}
    \label{fig:parallel}
    % \includegraphics[width=0.8\textwidth]{parallel_trend.pdf}
    \begin{minipage}{0.8\textwidth}\footnotesize
        注：图中横轴为相对时间（0 为政策实施年份），纵轴为 DID 系数，
        虚线表示 95\% 置信区间。
    \end{minipage}
\end{figure}

% ── 六、结论与启示 ────────────────────────────────────────────────────────
\section{结论与启示}
\label{sec:conclusion}

本文以\dots为研究对象，考察了\dots的影响及其机制。

本文的主要结论如下：第一，\dots； 第二，\dots。

本文的研究启示在于：第一，\dots；第二，\dots。

\textbf{参考文献：}
\bibliography{references}

% ── 附录 ──────────────────────────────────────────────────────────────────
\newpage
\appendix
\section{附录}

\subsection{稳健性检验}
\label{app:robust}

为保证研究结论的可靠性，本文进行了如下稳健性检验：

（1）\textbf{替换核心变量：}采用\dots作为\dots的替代变量，主要结论保持不变。

（2）\textbf{改变样本范围：}剔除\dots样本后重新回归，结果仍支持本文假说。

（3）\textbf{改变估计方法：}采用 Poisson 回归和 Tobit 模型，结果保持一致。

（4）\textbf{调整控制变量：}参考\citet{xxx}的经典文献，加入\dots等变量，结果稳健。

\end{document}
